\documentclass[report.tex]{subfiles}

\begin{document}

\section{The Register Transfer Language (RTL)}
\label{sec:rtl}

To evaluate our framework on realistic compiler intermediate
representations,
we formalize a low-level imperative language based on the
Register Transfer Language (RTL) of CompCert~\parencite{leroy-09-compcert}.
This language has two kinds of storage:
local register banks attached to execution states,
which act as local variables,
and a global shared heap memory.

\subsection{Values, Registers, and Memory}
\label{sec:rtl-val-mem}

Computations in RTL manipulate first-order values.
There are four kinds of values,
where $\loc$ represents allocated heap locations:
\begin{align*}
  v \in \val
  &
    \Coloneqq \VInt{z} \mid \VBool{b} \mid \VPtr{l} \mid \VUndef
  \\
  &
    \quad \text{with } z \in \mathbb{Z},
    b \in \{\true, \false\},
    \text{ and } l \in \loc
\end{align*}
$\VUndef$ represents an undefined value,
used to model uninitialized data.
It can be copied between registers,
stored into memory cells,
or passed as function arguments.
However,
attempting to execute an instruction that requires a specific kind of value,
such as an arithmetic calculation,
will cause the execution to get stuck.

To relate source and target values,
we define a similarity predicate $\sameval{I}{v_t}{v_s}$.
This predicate is parameterized by a binary relation on locations,
$I \subseteq \loc \times \loc$.
This relation will be made precise later in \cref{sec:mem-predicate}.
The deduction rules of the value similarity predicate are given in
\cref{fig:same-val}.
We note that the rule \textsc{Val-Undef} is asymmetric:
whenever the source program produces an undefined value $\VUndef$,
it can be related to any value $v_t$ produced by the target program.
This asymmetry allows the compiler to replace undefined source computations with
concrete target operations.

\begin{figure}[htp]
  \begin{equation*}
    \begin{array}{c@{\qquad}c}
      \begin{prooftree}
        \hypo{z_t = z_s}
        \infer1[\textsc{Val-Int}]{\sameval{I}{\VInt{z_t}}{\VInt{z_s}}}
      \end{prooftree}
      &
      \begin{prooftree}
        \hypo{b_t = b_s}
        \infer1[\textsc{Val-Bool}]{\sameval{I}{\VBool{b_t}}{\VBool{b_s}}}
      \end{prooftree}
      \\[2em]
      \begin{prooftree}
        \hypo{(l_t, l_s) \in I}
        \infer1[\textsc{Val-Ptr}]{\sameval{I}{\VPtr{l_t}}{\VPtr{l_s}}}
      \end{prooftree}
      &
      \begin{prooftree}
        \hypo{}
        \infer1[\textsc{Val-Undef}]{\sameval{I}{v_t}{\VUndef}}
      \end{prooftree}
    \end{array}
  \end{equation*}
  \caption{The value similarity predicate.}
  \label{fig:same-val}
\end{figure}

In RTL,
registers store values that can be actively manipulated in computations.
They are referenced by an identifier $r \in \reg$.
We do not restrict the number of registers.
An RTL program can also interact with the heap.
The heap is indexed by locations $l \in \loc$,
and the program can read,
write,
allocate,
and free memory cells associated with locations.

\subsection{Syntax and Program Structure}

High-level imperative languages employ \emph{structured} control-flow constructs,
such as nested conditionals,
while loops,
and compound statements.
In structured languages,
execution paths follow the hierarchical source code structure of the program.
In contrast,
RTL is an \emph{unstructured} language.
Control flow in RTL is represented as an explicit directed graph whose nodes
$pc \in \node$ represent program points.
Each instruction explicitly specifies the program points of its successors,
eliminating the need for loop or block constructs.
The instruction set and operator grammar of RTL are defined in
\cref{fig:rtl-grammar}.
An RTL function $f$ is a record containing:
\begin{equation*}
  f \in \rtlfun \Coloneqq \left\{\,
  \begin{array}{l@{\quad}l}
    \fnname \in \ident
    & \text{function identifier} \\
    \fnregs \in \listT{\reg}
    & \text{parameter list} \\
    \fnentry \in \node
    & \text{function entry point} \\
    \fncode \in \node \rightharpoonup \instr
    & \text{control-flow graph}
  \end{array}
  \,\right\}
\end{equation*}
We enforce the invariant that input parameter registers contain no duplicates.
The code map $\fncode$ is a finite map associating program points
with instructions.
In the following,
we write $\pcatis{f}{pc}{i}$ to assert that the code map of $f$ contains
instruction $i$ at program point $pc$.
An RTL program $P$ consists of a list of functions with a distinguished
entry function.
We require that function names in the program are unique,
and that the entry function is present.

\begin{figure}[htp]
  \centering
  \begin{tabular}{ccll}
    $i \in \instr$
    & $\Coloneqq$
    & $\Inop{pc}$
    & unconditional jump \\
    & $\mid$
    & $\Iop{dst}{\textsf{op}}{\vec{r}}{pc}$
    & pure arithmetic or logical operation \\
    & $\mid$
    & $\Iload{dst}{src}{pc}$
    & heap memory load \\
    & $\mid$
    & $\Istore{addr}{src}{pc}$
    & heap memory store \\
    & $\mid$
    & $\Ialloc{dst}{pc}$
    & heap cell allocation \\
    & $\mid$
    & $\Ifree{addr}{pc}$
    & heap cell deallocation \\
    & $\mid$
    & $\Icond{cond}{pc_{\true}}{pc_{\false}}$
    & conditional branching \\
    & $\mid$
    & $\Icall{dst}{fn}{\vec{r}}{pc}$
    & function invocation \\
    & $\mid$
    & $\Iret{src}$
    & function return
    \\[1em]
    $\textsf{op}$
    & $\Coloneqq$
    & $\Oadd \mid \Osub \mid \Omul \mid \Odiv$
    & integer arithmetic \\
    & $\mid$
    & $\Oincr \mid \Odecr$
    & increment and decrement \\
    & $\mid$
    & $\Oimmint{z} \mid \Oimmbool{b}$
    & constant immediate values \\
    & $\mid$
    & $\Omove \mid \Oeqz$
    & register copy and zero test
  \end{tabular}
  \\[0.8em]
  {\small
    where $dst, src, addr, cond \in \reg$,
    $\vec{r} \in \listT{\reg}$,
    $fn \in \ident$,
    and $pc, pc_{\true}, pc_{\false} \in \node$.
  }
  \caption{The syntax of RTL instructions and operators.}
  \label{fig:rtl-grammar}
\end{figure}

\subsection{Small-Step Operational Semantics}

The execution of an RTL program is formalized as a small-step operational
transition system.
The configurations are tuples $(\sigma, s, m)$,
where $\sigma$ is a call stack,
$m \in \mem$ is the global heap,
and $s \in \pcstate$ is a control state which can take three distinct forms:
\begin{itemize}
  \item $\State{f}{pc}{\rho}$ is the main state.
        It refers to an execution occurring inside function $f$
        at program point $pc$,
        with local register bank $\rho$.

  \item $\CallState{f}{\vec{v}}$
        represents the beginning of the invocation of function $f$
        with argument values $\vec{v}$.

  \item $\ReturnState{v}$ represents the completion of a function execution
        with return value $v$.
\end{itemize}

The call stack $\sigma$ is a list of stack frames.
Each stack frame $\Stackframe{dst}{f}{pc}{\rho}$ stores the caller
context:
the destination register $dst$ designated to receive the return value,
the caller function $f$,
the return program point $pc$,
and the caller register bank $\rho$.
Structuring transitions through explicit $\CallState{f}{\vec{v}}$ and
$\ReturnState{v}$ configurations cleanly decouples caller and callee
environments.
In CompCert~\parencite{leroy-09-compcert},
this abstraction is also used to integrate external function calls that jump
directly from a call state to a return state.

A register bank and the memory heap are formalized as shown in
\cref{fig:mem-and-reg}.
By explicitly recording $\Freed$ cells rather than removing locations
from the domain of $m$,
our memory model distinguishes dangling pointer dereferences and double frees
from accesses to unallocated memory.

\begin{figure}[htp]
  \centering
  \begin{equation*}
    \begin{array}{r@{\;\in\;}l@{\quad\Coloneqq\quad}l@{\qquad}l}
      \rho & \regbank & \reg \to \val & \text{register bank} \\
      m & \mem & \loc \rightharpoonup \cell & \text{memory heap} \\
      c & \cell & \Allocated{v} \mid \Freed & \text{heap memory cell}
    \end{array}
  \end{equation*}
  \vspace{0.5em}
  \begin{tabular}{ll}
    \toprule
    Notation & Operation Semantics \\
    \midrule
    $\getreg{r}{\rho}$ & Read register $r$;
    yields $\VUndef$ if $r \notin \dom{\rho}$ \\
    $\getreg{\vec{r}}{\rho}$ & Pointwise extension to register lists,
    returning $\vec{v}$ \\
    $\setreg{r}{v}{\rho}$ & Register assignment updating $r$ to value $v$ \\
    \addlinespace
    $\getat{l}{m}$ & Read cell at location $l$ \\
    $\setat{l}{v}{m}$ & Update allocated cell at $l$ to value $v$ \\
    $\allocat{l}{v}{m}$ & Allocate fresh cell $l \notin \dom{m}$
    initialized with value $v$ \\
    $\freeat{l}{m}$ & Deallocate active cell at $l$,
    setting its status to $\Freed$ \\
    \bottomrule
  \end{tabular}
  \caption{Formalization of register banks and memory heap operations.}
  \label{fig:mem-and-reg}
\end{figure}

\begin{figure}[htb]
  \begin{gather*}
    % \begin{prooftree}
    %   \hypo{\pcatis{f}{pc}{\Inop{pc'}}}
    %   \infer1[\textsc{Step-Nop}]{
    %     \step{P}{
    %       (\sigma, \State{f}{pc}{\rho}, m)
    %     }{
    %       (\sigma, \State{f}{pc'}{\rho}, m)
    %     }
    %   }
    % \end{prooftree}
    % \\[1em]
    \begin{prooftree}
      \hypo{\pcatis{f}{pc}{\Iop{dst}{op}{\vec{r}}{pc'}}}
      \hypo{\getreg{\vec{r}}{\rho} = \vec{v}}
      \hypo{\evalop{op}{\vec{v}} = \Some{v}}
      \infer3[\textsc{Step-Op}]{
        \step{P}{
          (\sigma, \State{f}{pc}{\rho}, m)
        }{
          (\sigma, \State{f}{pc'}{\setreg{dst}{v}{\rho}}, m)
        }
      }
    \end{prooftree}
    % \\[1em]
    % \begin{prooftree}
    %   \hypo{\pcatis{f}{pc}{\Iload{dst}{addr}{pc'}}}
    %   \hypo{\getreg{addr}{\rho} = \VPtr{l}}
    %   \hypo{\getat{l}{m} = \Allocated{v}}
    %   \infer3[\textsc{Step-Load}]{
    %     \step{P}{
    %       (\sigma, \State{f}{pc}{\rho}, m)
    %     }{
    %       (\sigma, \State{f}{pc'}{\setreg{dst}{v}{\rho}}, m)
    %     }
    %   }
    % \end{prooftree}
    % \\[1em]
    % \begin{prooftree}
    %   \hypo{\pcatis{f}{pc}{\Istore{addr}{src}{pc'}}}
    %   \hypo{\getreg{addr}{\rho} = \VPtr{l}}
    %   \hypo{\getreg{src}{\rho} = v}
    %   \hypo{\getat{l}{m} = \Allocated{\_}}
    %   \infer4[\textsc{Step-Store}]{
    %     \step{P}{
    %       (\sigma, \State{f}{pc}{\rho}, m)
    %     }{
    %       (\sigma, \State{f}{pc'}{\rho}, \setat{l}{v}{m})
    %     }
    %   }
    % \end{prooftree}
    \\[1em]
    \begin{prooftree}
      \hypo{\pcatis{f}{pc}{\Ialloc{dst}{pc'}}}
      \hypo{l \notin \dom{m}}
      \hypo{v \in \val}
      \infer3[\textsc{Step-Alloc}]{
        \step{P}{
          (\sigma, \State{f}{pc}{\rho}, m)
        }{
          (\sigma,
           \State{f}{pc'}{\setreg{dst}{\VPtr{l}}{\rho}},
           \allocat{l}{v}{m})
        }
      }
    \end{prooftree}
    \\[1em]
    \begin{prooftree}
      \hypo{\pcatis{f}{pc}{\Ifree{addr}{pc'}}}
      \hypo{\getreg{addr}{\rho} = \VPtr{l}}
      \hypo{\getat{l}{m} = \Allocated{\_}}
      \infer3[\textsc{Step-Free}]{
        \step{P}{
          (\sigma, \State{f}{pc}{\rho}, m)
        }{
          (\sigma, \State{f}{pc'}{\rho}, \freeat{l}{m})
        }
      }
    \end{prooftree}
    % \\[1em]
    % \begin{prooftree}
    %   \hypo{\pcatis{f}{pc}{\Icond{cond}{pc_{\true}}{pc_{\false}}}}
    %   \hypo{\getreg{cond}{\rho} = \VBool{b}}
    %   \infer2[\textsc{Step-Cond}]{
    %     \step{P}{
    %       (\sigma, \State{f}{pc}{\rho}, m)
    %     }{
    %       (\sigma, \State{f}{(b \;?\; pc_{\true} : pc_{\false})}{\rho}, m)
    %     }
    %   }
    % \end{prooftree}
    \\[1em]
    \begin{prooftree}
      \hypo{\pcatis{f}{pc}{\Icall{dst}{fn}{\vec{r}}{pc'}}}
      \hypo{\findfun{P}{fn} = \Some{f_{\text{callee}}}}
      \hypo{\getreg{\vec{r}}{\rho} = \vec{v}}
      \infer3[\textsc{Step-Call}]{
        \step{P}{
          (\sigma, \State{f}{pc}{\rho}, m)
        }{
          (\Stackframe{dst}{f}{pc'}{\rho} :: \sigma,
           \CallState{f_{\text{callee}}}{\vec{v}},
           m)
        }
      }
    \end{prooftree}
    \\[1em]
    \begin{prooftree}
      \hypo{|\vec{v}| = |\fnregs(f)|}
      \hypo{\rho = \initregs{\fnregs(f)}{\vec{v}}}
      \infer2[\textsc{Step-Entry}]{
        \step{P}{
          (\sigma, \CallState{f}{\vec{v}}, m)
        }{
          (\sigma, \State{f}{\fnentry(f)}{\rho}, m)
        }
      }
    \end{prooftree}
    \\[1em]
    \begin{prooftree}
      \hypo{\pcatis{f}{pc}{\Iret{src}}}
      \hypo{\getreg{src}{\rho} = v}
      \infer2[\textsc{Step-Ret}]{
        \step{P}{
          (\sigma, \State{f}{pc}{\rho}, m)
        }{
          (\sigma, \ReturnState{v}, m)
        }
      }
    \end{prooftree}
    \\[1em]
    \begin{prooftree}
      \infer0[\textsc{Step-Return}]{
        \step{P}{
          (\Stackframe{dst}{f}{pc'}{\rho} :: \sigma, \ReturnState{v}, m)
        }{
          (\sigma, \State{f}{pc'}{\setreg{dst}{v}{\rho}}, m)
        }
      }
    \end{prooftree}
  \end{gather*}
  \caption{Selected rules of the small-step operational semantics of RTL.}
  \label{fig:rtl-semantics}
\end{figure}

The step relation $\step{P}{s}{s'}$ is defined by inference rules,
of which are presented in \cref{fig:rtl-semantics}.
Execution of a function progresses through three main phases
($\textsf{CallState} \to \textsf{State} \to \textsf{ReturnState}$):
\begin{enumerate}
  \item In the beginning,
        the state is a call state $\CallState{f}{\vec{v}}$.
        We use the rule \textsc{Step-Entry}
        to step into the body of the function at its entry point.
        This initializes the register bank with the values given for the input
        registers of $f$.
        The resulting state is an active state $\State{f}{pc_0}{\rho}$.

  \item Then we use the instruction rules to execute the body of the function.
        These rules can only apply if they correspond to the next instruction
        to be executed,
        referenced by $pc$ in $f$.
        We note that evaluation of some instructions is partial.
        For instance,
        in the \textsc{Step-Op} rule,
        if $\evalop{op}{\vec{v}}$ returns $\None$, the execution is stuck.

  \item When the function reaches a return instruction,
        the rule \textsc{Step-Ret} is applied.
        It evaluates its return register and enters a return state
        $\ReturnState{v}$.

  \item If the call stack is not empty,
        the caller resumes execution via \textsc{Step-Return},
        popping the top stack frame,
        writing the return value into register $dst$,
        and branching to the next $pc$.

  \item When a function call occurs,
        the rule \textsc{Step-Call} is used.
        The caller evaluates the argument registers,
        pushes a fresh stack frame recording its continuation point $pc'$ and
        register bank $\rho$,
        and transitions to a call state $\CallState{f_{\text{callee}}}{\vec{v}}$.
\end{enumerate}

An execution state is final when the call stack is empty and the program is in
a return state. The final value is the one stored in the control state.

\subsection{Stack Modularity and Execution Invariants}

Reasoning directly about executions with arbitrary call stacks introduces
unnecessary proof complexity.
Without stack modularity,
one would need to maintain an explicit invariant over caller stack frames and
prove that it is preserved across every step.
To achieve modularity,
we establish lemmas showing that function executions can be analyzed in
isolation with an empty call stack $\sigma = []$,
and then embedded into arbitrary execution contexts.

% Operational steps are preserved when extending the call stack
% by an arbitrary frame list $\Sigma$:
% \begin{lemma}[Stack Lifting]
% \label{lem:rtl-lift-step}
% For all call stacks $\sigma, \sigma', \Sigma$,
% control states $ps, pt$,
% and memories $m, m'$:
% \begin{align*}
%   \step{P}{(\sigma, ps, m)}{(\sigma', pt, m')}
%   & \implies
%   \step{P}{(\sigma \mathbin{+\!+} \Sigma, ps, m)}{(\sigma' \mathbin{+\!+} \Sigma, pt, m')},
%   \\
%   \steps{P}{(\sigma, ps, m)}{(\sigma', pt, m')}
%   & \implies
%   \steps{P}{(\sigma \mathbin{+\!+} \Sigma, ps, m)}{(\sigma' \mathbin{+\!+} \Sigma, pt, m')}.
% \end{align*}
% \end{lemma}

% Conversely,
% the unlifting lemma shows that transitions performed in a populated
% stack context reflect reductions of the local frame:
% \begin{lemma}[Stack Unlifting]
% \label{lem:rtl-unlift-step}
% For all call stacks $\sigma, \Sigma$,
% control states $ps$,
% memories $m$,
% and target configurations $t''$:

% If $(\sigma, ps, m)$ is not a final value
% and $\step{P}{(\sigma \mathbin{+\!+} \Sigma, ps, m)}{t''}$,
% then there exist a local stack $\sigma'$,
% control state $pt$,
% and memory $m'$ such that:
% \begin{equation*}
%   t'' = (\sigma' \mathbin{+\!+} \Sigma, pt, m')
%   \quad \text{and} \quad
%   \step{P}{(\sigma, ps, m)}{(\sigma', pt, m')}.
% \end{equation*}
% \end{lemma}

% These lemmas provide the foundation for modular verification in our logic.
When verifying a function call,
the callee is verified starting from an empty call stack
$([], \CallState{f}{\vec{v}}, m)$.
Once the callee reduces to $([], \ReturnState{v}, m')$,
the derivation is lifted back into the caller's stack context,
reconnecting the callee's execution with the caller's continuation.

\end{document}

% LocalWords:  Coinductive coinductive asymmetricity arities
